\documentclass[11pt,a4paper]{article}
\usepackage[margin=1in]{geometry}
\usepackage{amsmath,amssymb,amsthm}
\usepackage{algorithm}
\usepackage{algpseudocode}
\usepackage{booktabs}
\usepackage{hyperref}

\theoremstyle{definition}
\newtheorem{definition}{Definition}[section]
\newtheorem{theorem}{Theorem}[section]
\newtheorem{lemma}[theorem]{Lemma}
\newtheorem{remark}{Remark}[section]

\title{\textbf{Theoretical Foundations, Invariance Characteristics, and Per-Example Statistical Mechanics of Layer Normalization}}
\author{Team Scaibu \\ \small Email: \texttt{jaydeep.9664920749@gmail.com} \\ \small Architecture and Core Engine Team}
\date{October 2026}

\begin{document}
\maketitle

\begin{abstract}
This document provides the complete mathematical specification, gradient dynamics, and invariance theory of Layer Normalization (LayerNorm). Unlike batch-dependent normalization schemes, LayerNorm computes sample statistics independently across the feature dimension of each individual token or example. We formulate the forward transformation, prove shift and scale invariance of pre-activations, detail the computational and memory bounds, step through a numerical trace completely consistent with the reference implementation, and analyze floating-point numerical considerations.
\end{abstract}

\section{Notation and Mathematical Preliminaries}

Let $X \in \mathbb{R}^{B \times D}$ represent an activation matrix where each row $x_i \in \mathbb{R}^D$ corresponds to an independent token or sample across $B$ items and $D$ hidden features.

\begin{itemize}
    \item $B \in \mathbb{N}_{\ge 1}$: Batch size or sequence length.
    \item $D \in \mathbb{N}_{\ge 1}$: Hidden feature dimension.
    \item $x_i \in \mathbb{R}^D$: Feature vector of sample $i \in \{1, \dots, B\}$.
    \item $\mu_i \in \mathbb{R}$: Sample mean computed over all $D$ channels for token $i$.
    \item $\sigma_i^2 \in \mathbb{R}_{\ge 0}$: Sample variance across all $D$ channels for token $i$.
    \item $\hat{x}_i \in \mathbb{R}^D$: Zero-mean, unit-variance normalized vector for sample $i$.
    \item $\gamma \in \mathbb{R}^D$: Channel-wise learned scale parameter vector.
    \item $\beta \in \mathbb{R}^D$: Channel-wise learned bias parameter vector.
    \item $\epsilon > 0$: Numerical denominator regularizer ($\epsilon = 10^{-5}$).
\end{itemize}

\begin{table}[h]
\centering
\caption{Summary of Mathematical Symbols for Layer Normalization}
\begin{tabular}{lll}
\toprule
\textbf{Symbol} & \textbf{Type / Domain} & \textbf{Description} \\
\midrule
$X$ & $\mathbb{R}^{B \times D}$ & Activation input matrix \\
$x_i$ & $\mathbb{R}^D$ & Activation vector of example $i$ \\
$\mu$ & $\mathbb{R}^B$ & Vector of per-example means \\
$\sigma^2$ & $\mathbb{R}_{\ge 0}^B$ & Vector of per-example variances \\
$\hat{X}$ & $\mathbb{R}^{B \times D}$ & Standardized activation tensor \\
$Y$ & $\mathbb{R}^{B \times D}$ & Layer normalized output tensor \\
$\gamma$ & $\mathbb{R}^D$ & Gain parameter vector \\
$\beta$ & $\mathbb{R}^D$ & Bias parameter vector \\
$\epsilon$ & $\mathbb{R}_{> 0}$ & Numerical stability constant \\
\bottomrule
\end{tabular}
\end{table}

\section{Problem Definition and Core Concepts}

\subsection{Intuitive Meaning}
While Batch Normalization pools statistics along the batch axis (introducing cross-sample coupling and failing on small batch sizes or variable-length recurrent sequences), Layer Normalization calculates mean and variance strictly across the feature dimensions of a single example. This enables identical computation during training and inference, regardless of batch size.

\subsection{Formal Mathematical Definition}
For each sample $i \in \{1, \dots, B\}$:
\begin{align}
\mu_i &= \frac{1}{D} \sum_{j=1}^D X_{ij} \\
\sigma_i^2 &= \frac{1}{D} \sum_{j=1}^D (X_{ij} - \mu_i)^2 \\
\hat{X}_{ij} &= \frac{X_{ij} - \mu_i}{\sqrt{\sigma_i^2 + \epsilon}} \\
Y_{ij} &= \gamma_j \hat{X}_{ij} + \beta_j
\end{align}

\section{Algorithm Specification}

\begin{algorithm}[H]
\caption{Layer Normalization Forward Transformation}
\label{alg:layer_norm}
\begin{algorithmic}[1]
\Require Input activation matrix $X \in \mathbb{R}^{B \times D}$, gain vector $\gamma \in \mathbb{R}^D$, bias vector $\beta \in \mathbb{R}^D$, regularizer $\epsilon > 0$.
\Ensure Normalized output matrix $Y \in \mathbb{R}^{B \times D}$, mean vector $\mu \in \mathbb{R}^B$, variance vector $\sigma^2 \in \mathbb{R}^B$.
\State $B \gets \text{rows}(X)$, $D \gets \text{cols}(X)$
\State Allocate $Y \in \mathbb{R}^{B \times D}$, $\mu \in \mathbb{R}^B$, $\sigma^2 \in \mathbb{R}^B$
\For{$i \gets 1 \textbf{ to } B$}
    \State $\mu_i \gets \frac{1}{D} \sum_{j=1}^D X_{ij}$
    \State $\sigma_i^2 \gets \frac{1}{D} \sum_{j=1}^D (X_{ij} - \mu_i)^2$
    \State $s_i \gets \sqrt{\sigma_i^2 + \epsilon}$
    \For{$j \gets 1 \textbf{ to } D$}
        \State $\hat{X}_{ij} \gets \frac{X_{ij} - \mu_i}{s_i}$
        \State $Y_{ij} \gets \gamma_j \hat{X}_{ij} + \beta_j$
    \EndFor
\EndFor
\State \Return $\{Y, \mu, \sigma^2\}$
\end{algorithmic}
\end{algorithm}

\section{Theoretical Guarantees and Invariance Properties}

\begin{theorem}[Shift and Scaling Invariance of Pre-activations]
Let $x \in \mathbb{R}^D$. For any scalar shift $c \in \mathbb{R}$ and scalar multiplier $\alpha > 0$, the standardized representation is invariant:
\begin{equation}
\text{LN}(\alpha x + c \mathbf{1}) = \text{LN}(x)
\end{equation}
\end{theorem}

\begin{proof}
Let $x' = \alpha x + c \mathbf{1}$. The transformed mean is $\mu' = \frac{1}{D}\sum_{j=1}^D (\alpha x_j + c) = \alpha \mu + c$.
The transformed variance is:
\begin{equation}
\sigma'^2 = \frac{1}{D}\sum_{j=1}^D ((\alpha x_j + c) - (\alpha \mu + c))^2 = \frac{1}{D}\sum_{j=1}^D \alpha^2 (x_j - \mu)^2 = \alpha^2 \sigma^2
\end{equation}
Substituting into the normalization formula (neglecting $\epsilon$):
\begin{equation}
\hat{x}'_j = \frac{(\alpha x_j + c) - (\alpha \mu + c)}{\sqrt{\alpha^2 \sigma^2}} = \frac{\alpha (x_j - \mu)}{\alpha \sigma} = \frac{x_j - \mu}{\sigma} = \hat{x}_j
\end{equation}
Applying the element-wise transformation $Y_j = \gamma_j \hat{x}_j + \beta_j$ yields identical outputs.
\end{proof}

\subsection{Limitations}
\begin{itemize}
    \item \textbf{High Spatial Variance in CNNs}: When feature dimensions vary widely in semantic meaning (e.g., across edge detectors vs. color channels in conv layers), normalizing across all channels can dampen informative channel-specific scale variations.
\end{itemize}

\section{Complexity Analysis}

\subsection{Time Complexity}
\begin{itemize}
    \item \textbf{Per-row Reductions}: Computing row mean and row variance requires $2D$ operations per sample.
    \item \textbf{Per-element Scaling}: Standardizing and applying $\gamma, \beta$ requires $3D$ operations per sample.
    \item \textbf{Total Time Complexity}: $\Theta(B \cdot D)$ FLOPs.
\end{itemize}

\subsection{Space Complexity}
\begin{itemize}
    \item \textbf{Output Matrix}: Requires storing $B \times D$ elements.
    \item \textbf{Auxiliary Memory}: Requires $\mu \in \mathbb{R}^B$ and $\sigma^2 \in \mathbb{R}^B$, yielding $\mathcal{O}(B)$ auxiliary storage.
\end{itemize}

\section{Exactness and Error Analysis}

Layer normalization is exact up to floating-point precision. Summing across hidden dimension $D$ using double-precision accumulators ensures numerical stability and prevents catastrophic cancellation.

\section{Worked Numerical Example}

Let $B = 2$ tokens with hidden dimension $D = 3$:
\begin{equation}
X = \begin{bmatrix} 1.0 & 2.0 & 3.0 \\ 2.0 & 4.0 & 6.0 \end{bmatrix}, \quad \gamma = [1.0, 1.0, 1.0]^T, \quad \beta = [0.0, 0.0, 0.0]^T, \quad \epsilon = 10^{-5}
\end{equation}

\begin{enumerate}
    \item \textbf{Row 1 Statistics}:
    \begin{align}
    \mu_1 &= \frac{1.0 + 2.0 + 3.0}{3} = 2.0 \\
    \sigma_1^2 &= \frac{(1-2)^2 + (2-2)^2 + (3-2)^2}{3} = \frac{1 + 0 + 1}{3} = \frac{2}{3} \approx 0.6666667
    \end{align}
    Standard deviation: $\sqrt{0.6666667 + 10^{-5}} = \sqrt{0.6666767} \approx 0.8165027$.
    \item \textbf{Row 1 Normalized Values}:
    \begin{align}
    \hat{X}_{11} &= \frac{1.0 - 2.0}{0.8165027} \approx -1.224736, \quad Y_{11} \approx -1.224736 \\
    \hat{X}_{12} &= \frac{2.0 - 2.0}{0.8165027} = 0.0, \quad Y_{12} = 0.0 \\
    \hat{X}_{13} &= \frac{3.0 - 2.0}{0.8165027} \approx 1.224736, \quad Y_{13} \approx 1.224736
    \end{align}
    \item \textbf{Row 2 Statistics}:
    \begin{align}
    \mu_2 &= \frac{2.0 + 4.0 + 6.0}{3} = 4.0 \\
    \sigma_2^2 &= \frac{(2-4)^2 + (4-4)^2 + (6-4)^2}{3} = \frac{4 + 0 + 4}{3} = \frac{8}{3} \approx 2.6666667
    \end{align}
    Standard deviation: $\sqrt{2.6666667 + 10^{-5}} = \sqrt{2.6666767} \approx 1.6330024$.
    \item \textbf{Row 2 Normalized Values}:
    \begin{align}
    \hat{X}_{21} &= \frac{2.0 - 4.0}{1.6330024} \approx -1.224736, \quad Y_{21} \approx -1.224736 \\
    \hat{X}_{22} &= \frac{4.0 - 4.0}{1.6330024} = 0.0, \quad Y_{22} = 0.0 \\
    \hat{X}_{23} &= \frac{6.0 - 4.0}{1.6330024} \approx 1.224736, \quad Y_{23} \approx 1.224736
    \end{align}
\end{enumerate}

\section{Numerical Considerations and Edge Cases}

\begin{itemize}
    \item \textbf{Flat Activations}: When all features in a vector are constant ($x_{ij} = c$), variance $\sigma_i^2 = 0$. Epsilon $\epsilon > 0$ guarantees finite division.
    \item \textbf{Mixed-Precision Arithmetic}: In FP16/BF16 training, calculating sum of squares $\sum X_{ij}^2$ can overflow or underflow. Accumulation should occur in FP32 before casting back.
\end{itemize}

\begin{thebibliography}{9}
\bibitem{ba2016layer} J.~L.~Ba, J.~R.~Kiros, and G.~E.~Hinton, ``Layer Normalization,'' \emph{arXiv preprint arXiv:1607.06450}, 2016.
\bibitem{vaswani2017attention} A.~Vaswani et al., ``Attention Is All You Need,'' in \emph{Advances in Neural Information Processing Systems (NeurIPS)}, pp.~5998--6008, 2017.
\end{thebibliography}

\end{document}
